\documentclass[10pt,a4paper]{amsart}
\usepackage[margin=19mm]{geometry}
\usepackage{amsmath,amssymb,mathtools}
\usepackage[scr=rsfs,cal=euler]{mathalfa}
\usepackage{xfrac}
\usepackage[unicode,hidelinks,bookmarks=false]{hyperref}
\newcommand{\ie}{i.e.\ }
\newcommand{\cf}{cf.\ }
\newcommand{\resp}{respectively\ }
\newcommand{\Subsection}{\subsection}
\providecommand{\nobreakdash}{\nobreak\hbox{-}}
\setlength{\emergencystretch}{2em}
\multlinegap0pt
\usepackage{amssymb}
\usepackage{mathtools}
\usepackage{enumerate}
\usepackage{tikz}
\usepackage[normalem]{ulem}
\newtheorem{lemma}{Lemma}
\title{Proof of a Conjecture of Cui, Gu and Tang on 18-Colored Generalized Frobenius Partitions}
\author{Hirakjyoti Das, Hemjyoti Nath, Abhishek Sarma}
\date{Source: arXiv:2609.05302v1 (CC BY 4.0)}
\begin{document}
\maketitle

\noindent\emph{Source and licence.} Proof of a Conjecture of Cui, Gu and Tang on 18-Colored Generalized Frobenius Partitions. Version arXiv:2609.05302v1 (CC BY 4.0), distributed by arXiv under the Creative Commons Attribution 4.0 International licence, http://creativecommons.org/licenses/by/4.0/, as recorded in the arXiv licence metadata for this exact version and shown on the abstract page https://arxiv.org/abs/2609.05302v1. Full licence notice: \texttt{licenses/arxiv-2609.05302-CC-BY-4.0.txt}.\par\smallskip

\noindent\textit{Editorial scope.} Counted unit: the article's lemma lem:main-pk-cong (source0.tex lines 1150--1155). The article's complete original proof of it is delivered with it (source0.tex lines 1157--1339, one \texttt{proof} environment of 183 lines). No mathematics is written, repaired or re-derived: the statement and the proof are the article's own lines, in the article's own notation, and no retained line is substituted or re-flowed.  Retained uncounted as the source-local context the proof needs: lines 767--808; lines 1024--1043.  Nothing was omitted from any retained span, so the package carries no omission record.  No retained line contains a citation, so the wrapper carries no bibliography and no bibliography entry.  No retained line contains a figure environment, an image-inclusion command or drawing code, so no image is delivered and the package declares no asset dependency.  Editorial formatting only, in addition: an AMS \texttt{amsart} preamble that restates the article's own declarations and macros used by the retained lines, namely the environments \texttt{lemma} on separate counters, together with the standard packages those lines need.  The retained environments print in this document as Lemma 1 in order of appearance, whereas the article prints the same items with the numbering of its own full document.  The complete native source of the article as distributed in the arXiv e-print for this exact version is bundled unchanged as \texttt{sources/209424/source0.tex}.
\begin{lemma}\label{lem:H-decomp}
We have
\begin{equation*}
H(q)=M\left(13+RA(p)+R^2B(p)\right),
\end{equation*}
where
\begin{align*}
A(p)
={}&
\frac{2p+1}{p-1}
-
\frac{12p(p+1)^2}
     {(p-1)^3(p+2)}
-
\frac{21p^2(p+1)^2}
     {4(p-1)^3(2p+1)}
+
\frac{144p(p+1)^4}
     {(p+2)^5(2p+1)^2}
\end{align*}
and
\begin{align*}
B(p)
={}&
\frac{p}{2(2p+1)^2}
-
\frac{3p^2}
     {4(p-1)^2(2p+1)}
+
\frac{12}{p-1}
+
\frac{6p(p+1)}
     {(p-1)^2(p+2)(2p+1)}
\\
&-
\frac{9p^2(p+1)^2}
     {2(p-1)^4(p+2)(2p+1)^2}
+
\frac{9p^3}
     {8(p-1)^2(2p+1)^2}.
\end{align*}
\end{lemma}

\begin{lemma}\label{lem:R-mod27}
Let
\begin{align*}
\rho(p)
&=
a+\frac{3p}{2a}
-\frac{9p^2(2a^2+p)}{8a^5},\\
Q(p)
&=
8p^6-36p^5+54p^4-61p^3
+54p^2-36p+8.
\end{align*}
Then
\begin{align*}
R&\equiv\rho(p)\pmod{27},\\
\rho(p)
&=
-\frac{Q(p)}{8(p-1)^5}.
\end{align*}
\end{lemma}

\begin{lemma}\label{lem:main-pk-cong}
We have
\begin{equation*}
13+RA(p)+R^2B(p)\equiv0\pmod{27},
\end{equation*}where $A(p)$ and $B(p)$ are defined in Lemma \ref{lem:H-decomp}.
\end{lemma}

\begin{proof}
Using Lemma \ref{lem:R-mod27}, we have
\begin{equation*}
RA(p)\equiv\rho(p)A(p)\pmod{27}
\end{equation*}
and
\begin{equation*}
R^2B(p)\equiv\rho(p)^2B(p)\pmod{27}.
\end{equation*}
Therefore
\begin{equation*}
13+RA(p)+R^2B(p)
\equiv
13+\rho(p)A(p)+\rho(p)^2B(p)
\pmod{27}.
\end{equation*}

We set
\begin{equation*}
A(p):=-\frac{P_1(p)}{4(p-1)^3(p+2)^5(2p+1)^2},
\end{equation*}
where
\begin{align*}
P_1(p)
={}&
10p^{10}
+413p^9
+3270p^8
+16161p^7
+42570p^6
\\
&+
60996p^5
+49224p^4
+22512p^3
+5664p^2
+512p
-128,
\end{align*}
and
\begin{equation*}
B(p):=\frac{P_2(p)}{8(p-1)^4(p+2)(2p+1)^2},
\end{equation*}
where
\begin{equation*}
P_2(p)
=
385p^6
+82p^5
-1523p^4
+500p^3
+836p^2
-232p
-192.
\end{equation*}
Therefore, using the last lemma, we obtain
\begin{align*}
13+\rho A+\rho^2B
={}&
13
+
\frac{Q(p)P_1(p)}
{32(p-1)^8(p+2)^5(2p+1)^2}
\\
&+
\frac{Q(p)^2P_2(p)}
{512(p-1)^{14}(p+2)(2p+1)^2}.
\end{align*}

We take $D(p):=512(p-1)^{14}(p+2)^5(2p+1)^2$, then
\begin{align*}
D(p)\left(13+\rho A+\rho^2B\right)
&=
6656(p-1)^{14}(p+2)^5(2p+1)^2
\\
&\quad
+
16(p-1)^6Q(p)P_1(p)
+
(p+2)^4Q(p)^2P_2(p)
\\
&=:N(p).
\end{align*}
Simplifying,
\iffalse
\begin{align*}
N(p)
={}&
25920p^{22}
+46656p^{21}
-607824p^{20}
+38880p^{19}
+1670220p^{18}
\\
&+
6375564p^{17}
-18091755p^{16}
-944190p^{15}
+48815865p^{14}
\\
&-
67331304p^{13}
-903096p^{12}
+125516088p^{11}
-153157824p^{10}
\\
&-
8136288p^9
+173945664p^8
-137396736p^7
-2951424p^6
\\
&+
53706240p^5
-22339584p^4
-884736p^3
+1880064p^2
-221184p.
\end{align*}

Hence
\fi
\begin{align*}
N(p)
=
27p\Bigl(
&960p^{21}
+1728p^{20}
-22512p^{19}
+1440p^{18}
+61860p^{17}
\\
&+
236132p^{16}
-670065p^{15}
-34970p^{14}
+1807995p^{13}
\\
&-
2493752p^{12}
-33448p^{11}
+4648744p^{10}
-5672512p^9
\\
&-
301344p^8
+6442432p^7
-5088768p^6
-109312p^5
\\
&+
1989120p^4
-827392p^3
-32768p^2
+69632p
-8192
\Bigr).
\end{align*}
Therefore
\begin{equation*}
D(p)\left(13+\rho A+\rho^2B\right)
\equiv0\pmod{27}.
\end{equation*}

Observe that
\begin{equation*}
D(0)
=
512(-1)^{14}(2)^5
=
2^{14},
\end{equation*}
which gives
\begin{equation*}
3\nmid D(0).
\end{equation*}
Therefore, $D(p)$ is a unit in $\mathbb{Z}_3[[p]]$. Thus
\begin{equation*}
13+\rho A+\rho^2B
\equiv0\pmod{27},
\end{equation*}
which proves the lemma.
\end{proof}

\end{document}
